\section*{Bab \ref{ch:b2:magma-emplacement} --- Magma dan Penempatannya}

\begin{solution}{exo:b2:magma-emplacement:1}
$T = \qty{1073}{K}$. Untuk $w = 1$, $\ln w = 0$: $\log_{10}\eta = -3.545 +
9601/(1073 - 195.7) = -3.545 + 10.94 = 7.40$. Untuk $w = 4$, $\ln 4 = 1.386$:
$A = -3.545 + 1.155 = -2.390$, $B = 9601 - 3283 = 6318$, $C = 195.7 + 44.7 =
240.4$, dan $\log_{10}\eta = -2.390 + 6318/832.6 = -2.390 + 7.588 = 5.20$.
Viskositasnya turun dengan faktor $10^{7.40 - 5.20} = 10^{2.20} \approx 160$.
\end{solution}

\begin{solution}{exo:b2:magma-emplacement:2}
Riolit: $p_s = (6/0.44)^2 = \qty{186}{MPa}$, pada kedalaman $\num{186e6}/(2600\times9.81)
= \qty{7.3}{km}$. Basal: $p_s = (0.5/0.34)^2 = \qty{2.16}{MPa}$, pada
$\num{2.16e6}/(2600\times9.81) = \qty{85}{m}$. Riolit mulai membentuk gelembung
jauh di dalam kerak, basal baru di dekat permukaan (jika tidak membawa CO$_2$).
\end{solution}

\begin{solution}{exo:b2:magma-emplacement:3}
$(1150 + 50)/2 = \qty{600}{\celsius}$ dan $(1150 + 400)/2 = \qty{775}{\celsius}$.
\omterm{def:b2:magma-emplacement:bodies}{Retas} di dalam batuan dingin terdinginkan lebih kuat pada kontaknya, dan memiliki
\omterm{def:b2:magma-emplacement:aureole}{tepi dingin} yang lebih lebar.
\end{solution}

\begin{solution}{exo:b2:magma-emplacement:4}
(a) \omterm{def:b2:magma-emplacement:bodies}{Sill}, sejajar dengan perlapisan. (b) \omterm{def:b2:magma-emplacement:bodies}{Retas}, memotong perlapisan.
(c) \omterm{def:b2:magma-emplacement:bodies}{Batolit}, lebih dari \qty{100}{km^2} dan dirakit dari banyak intrusi.
(d) \omterm{def:b2:magma-emplacement:bodies}{Lakolit}, yang telah mengubah atapnya menjadi kubah. (e) \omterm{def:b2:magma-emplacement:bodies}{Retas}: granit tidak berlapis,
sehingga lembaran apa pun di dalamnya adalah \omterm{def:b2:magma-emplacement:bodies}{retas}, apa pun orientasinya.
\end{solution}

\begin{solution}{exo:b2:magma-emplacement:5}
$(1 - 1.35\phi)^{-2.5}$: untuk $\phi = 0.05$, $0.9325^{-2.5} = 1.19$, sehingga
\qty{119}{Pa.s}, dibandingkan $1 + 2.5\times0.05 = 1.125$, \qty{113}{Pa.s}, dengan rumus
Einstein (dekat, sebagaimana mestinya pada batas encer). Untuk $0.20$: $0.73^{-2.5} =
2.20$, \qty{220}{Pa.s}; untuk $0.40$: $0.46^{-2.5} = 6.97$, \qty{697}{Pa.s}; untuk $0.60$:
$0.19^{-2.5} = 63.6$, sekitar \qty{6360}{Pa.s}. Kemacetan pada $\phi = 1/1.35 = 0.74$.
\end{solution}

\begin{solution}{exo:b2:magma-emplacement:6}
$G = 300\times9.81 = \qty{2943}{Pa/m}$. (a) $\bar u = 2943\times1^2/(12\times100) =
\qty{2.45}{m/s}$; $q = \qty{2.45}{m^2/s}$; $Re = 2700\times2.45\times1/100 = 66$, laminar;
$t = 10000/2.45 = \qty{4080}{s}$, sekitar \qty{1.1}{h}. (b) $\bar u =
2943/(12\times\num{e6}) = \qty{2.45e-4}{m/s}$, $t = \qty{4.08e7}{s} = 472$ hari; dengan
$w = \qty{10}{m}$, $\bar u$ seratus kali lebih besar, \qty{0.0245}{m/s}, dan $t =
\qty{4.08e5}{s} = 4.7$ hari. (c) \omterm{def:b2:magma-emplacement:bodies}{Retas} \qty{1}{m} mendingin dalam sekitar
$0.5^2/\num{e-6} = \qty{2.5e5}{s}$, yaitu 2.9 hari; \omterm{def:b2:magma-emplacement:bodies}{retas} \qty{10}{m} dalam $5^2/\num{e-6} =
\qty{2.5e7}{s}$, yaitu 289 hari. Basal (\qty{1.1}{h}) dan riolit dalam \omterm{def:b2:magma-emplacement:bodies}{retas}
yang lebar (4.7 hari) tiba sebelum \omterm{def:b2:magma-emplacement:bodies}{retasnya} membeku; riolit dalam
\omterm{def:b2:magma-emplacement:bodies}{retas} \qty{1}{m} (472 hari) tidak dapat.
\end{solution}

\begin{solution}{exo:b2:magma-emplacement:7}
Per kilogram \omterm{def:b2:magma-emplacement:magma}{magma}, $n = 10/18.015 = \qty{0.555}{mol}$ air dan $nRT =
0.555\times8.314\times1173 = \qty{5414}{J}$; lelehan yang tersisa menempati $0.99/2400 =
\qty{4.13e-4}{m^3}$. Pada \qty{100}{MPa}: $V = \qty{5.41e-5}{m^3}$, fraksi
$5.41/(5.41 + 41.3) = 0.12$. Pada \qty{10}{MPa}: $V = \qty{5.41e-4}{m^3}$, fraksi
$0.57$. Pada \qty{0.1}{MPa}: $V = \qty{0.0541}{m^3}$, fraksi $0.992$. Nilai pada
\qty{100}{MPa} paling tidak dapat diandalkan: uap air pada tekanan setinggi itu rapat dan
jauh dari ideal (dan sebagian besar airnya sebenarnya masih terlarut).
\end{solution}

\begin{solution}{exo:b2:magma-emplacement:8}
(a) $2300 + 60z = 2650$ memberikan $z = \qty{5.8}{km}$. (b) Lelehan itu memiliki
\qty{2580}{kg/m^3}: $z = 280/60 = \qty{4.7}{km}$, sekitar \qty{1.2}{km} lebih tinggi. (c)
Pada $z = \qty{7.8}{km}$ batuan memiliki $2300 + 60\times7.8 = \qty{2770}{kg/m^3}$;
gayanya $(2770 - 2650)\times9.81 = \qty{1.2}{kN}$ per meter kubik, ke atas.
\end{solution}

\begin{solution}{exo:b2:magma-emplacement:9}
$v = 2\times300\times9.81\times(1000)^2/(9\times\num{e19}) = \qty{6.5e-11}{m/s}$, atau
$\num{6.5e-11}\times\num{3.156e7} = \qty{2.1e-3}{m}$ per tahun, \qty{2.1}{mm/yr}. Naik
\qty{10}{km} memerlukan $\num{e4}/\num{2.1e-3} = \qty{4.8}{Ma}$. Skala waktu pendinginan
adalah $(1000)^2/\num{e-6} = \qty{e12}{s} = \qty{32}{ka}$: tubuh itu membeku setelah
naik sekitar $32000\times0.0021 \approx \qty{70}{m}$. \omterm{def:b2:magma-emplacement:ascent}{Diapir} magmatik tidak dapat melintasi
kerak semacam itu; \omterm{def:b2:magma-emplacement:ascent}{diapir} memerlukan sekeliling yang lebih panas dan lebih lunak.
\end{solution}

\begin{solution}{exo:b2:magma-emplacement:10}
(a) $x_{\mathrm{OH}} = \sqrt{K_2x_m} = \sqrt{K_1K_2f}$, sehingga $x = K_1f + \frac12\sqrt{K_1K_2f}$.
Perbandingan suku pertama terhadap suku kedua adalah $2\sqrt{K_1f/K_2}$: perbandingan itu mendekati nol
ketika $f \to 0$ (akar kuadrat mendominasi) dan tumbuh tanpa batas bersama $f$ (suku
linear mendominasi). (b) $K_1 f = \frac12\sqrt{K_1K_2 f}$ memberikan $K_1^2f^2 =
K_1K_2f/4$, sehingga $f = K_2/(4K_1)$. (c) $f = 0.2/(\num{4e-4}) = \qty{500}{MPa}$. Pada
\qty{100}{MPa}: $x_m = 0.0100$, $\frac12\sqrt{\num{e-4}\times0.2\times100} =
\frac12\times0.0447 = 0.0224$, sehingga $x = 0.0324$ dan bagian hidroksilnya
$0.0224/0.0324 = \qty{69}{\percent}$.
\end{solution}

\begin{solution}{exo:b2:magma-emplacement:11}
(a) $1000/\num{4e6} = \qty{2.5e-4}{km^3}$ per tahun, yaitu $\qty{2.5e5}{m^3}$ per tahun,
atau $\num{2.5e5}/\num{3.156e7} = \qty{0.0079}{m^3/s}$, sekitar delapan liter per sekon.
(b) \omterm{def:b2:magma-emplacement:bodies}{Retas} itu membawa $1000\times2\times0.8 = \qty{1600}{m^3/s}$, dan mengantarkan
$\qty{e12}{m^3}$ dalam \qty{6.25e8}{s}, sekitar 20 tahun: $\num{5e-6}$ dari
\qty{4}{Ma}. (c) Setiap lembaran memuat $20\times\num{e4}\times\num{e4} = \qty{2e9}{m^3} =
\qty{2}{km^3}$, sehingga 500 lembaran, satu setiap $\num{4e6}/500 = 8000$ tahun rata-rata.
(d) $t = (10)^2/(4\times0.65^2\times\num{0.7e-6}) = \qty{8.5e7}{s}$, yaitu 2.7 tahun. Setiap
lembaran membeku ribuan tahun sebelum lembaran berikutnya tiba: kecuali pasokannya
datang dalam semburan yang cepat, \omterm{def:b2:magma-emplacement:bodies}{pluton} itu tidak pernah menjadi satu tubuh \omterm{def:b2:magma-emplacement:magma}{magma} yang besar, seperti
diisyaratkan penentuan umur \omterm{def:b2:magma-emplacement:bodies}{pluton-pluton} besar.
\end{solution}

\begin{solution}{exo:b2:magma-emplacement:12}
(a) Pada $x$ tetap, $\ln(T - T_0) = \text{konstanta} - \frac12\ln t - x^2/(4\kappa t)$, yang
turunannya $-1/(2t) + x^2/(4\kappa t^2)$ lenyap pada $t^* = x^2/(2\kappa)$. Maka
$e^{-x^2/(4\kappa t^*)} = e^{-1/2}$ dan $\sqrt{4\pi\kappa t^*} = x\sqrt{2\pi}$, sehingga $\Delta
T_{\max} = (T_m - T_0)\,b\,e^{-1/2}/(x\sqrt{2\pi}) = (T_m - T_0)\,b/(x\sqrt{2\pi e})$. (b)
$\sqrt{2\pi e} = 4.13$: $x = 1100\times2/(100\times4.13) = \qty{5.3}{m}$, dicapai setelah
$t^* = 5.3^2/(2\times\num{0.7e-6}) = \qty{2.0e7}{s}$, sekitar 230 hari (pada jarak hanya 2.7
kali lebar \omterm{def:b2:magma-emplacement:bodies}{retas}, pendekatan sumber bidang itu kasar). (c) Pemanasannya berkurang
sebagai $b/x$: \omterm{def:b2:magma-emplacement:bodies}{retas} tipis membawa sedikit panas, dan batuan hanya terpanaskan secara berarti
dalam jarak beberapa kali lebar \omterm{def:b2:magma-emplacement:bodies}{retas}.
\end{solution}

\begin{solution}{pb:b2:magma-emplacement:1}
\textbf{1.} $\log_{10}\eta = -4.55 + 5510/(1423 - 611) = -4.55 + 6.786 = 2.24$;
$\eta = \qty{172}{Pa.s}$.

\textbf{2.} $-4.55 + 5510/762 = -4.55 + 7.231 = 2.68$, $\eta = \qty{480}{Pa.s}$:
faktor $480/172 = 2.8$ untuk \qty{50}{K}.

\textbf{3.} $172\times(1 - 1.35\times0.20)^{-2.5} = 172\times0.73^{-2.5} = 172\times2.20 =
\qty{378}{Pa.s}$.

\textbf{4.} $\phi = 1/1.35 = 0.74$. Bukti lapangan meletakkan akhir aliran secara keseluruhan
pada fraksi lelehan \qtyrange{30}{50}{\percent}, yaitu
\qtyrange{50}{70}{\percent} kristal: lebih awal, karena kristal nyata bukan
bola berukuran seragam.

\textbf{5.} Batuan dasar: $2800 - 2710 = \qty{90}{kg/m^3}$, \omterm{def:b2:magma-emplacement:magma}{magma} lebih ringan.
Sedimen: $2450 - 2710 = \qty{-260}{kg/m^3}$, \omterm{def:b2:magma-emplacement:magma}{magma} lebih rapat.

\textbf{6.} $p = 2450\times9.81\times3000 = \qty{72.1}{MPa}$.

\textbf{7.} $90\times9.81 = \qty{883}{Pa/m}$.

\textbf{8.} $\bar u = 883\times2.0^2/(12\times378) = 3532/4536 = \qty{0.78}{m/s}$.

\textbf{9.} $q = 0.78\times2.0 = \qty{1.56}{m^2/s}$; sepanjang \qty{1.0}{km}, sekitar
\qty{1560}{m^3/s}.

\textbf{10.} $Re = 2710\times0.78\times2.0/378 = 11$: laminar.

\textbf{11.} $12000/0.78 = \qty{15400}{s}$, sekitar \qty{4.3}{h}.

\textbf{12.} $(1.0)^2/\num{0.7e-6} = \qty{1.43e6}{s}$, yaitu 16.5 hari, jauh lebih lama daripada
\qty{4.3}{h}: \omterm{def:b2:magma-emplacement:bodies}{retas} tetap terbuka selama \omterm{def:b2:magma-emplacement:magma}{magma} mengalir.

\textbf{13.} Pada \qty{3.0}{km} \omterm{def:b2:magma-emplacement:magma}{magma} bertemu sedimen yang lebih ringan daripada dirinya:
dasar sedimen adalah \omterm{def:b2:magma-emplacement:ascent}{tingkat daya apung netralnya}. Di atasnya \omterm{def:b2:magma-emplacement:magma}{magma} akan
lebih rapat daripada sekelilingnya, sehingga \omterm{def:b2:magma-emplacement:magma}{magma} menyebar ke samping di sepanjang kontak sebagai
\omterm{def:b2:magma-emplacement:bodies}{sill}.

\textbf{14.} $p_s = (0.40/0.34)^2 = \qty{1.38}{MPa}$, yang dicapai pada
$\num{1.38e6}/(2450\times9.81) = \qty{58}{m}$ di bawah permukaan.

\textbf{15.} $0.34\times\sqrt{72.1} = 0.34\times8.49 = \qty{2.9}{\percent}$, jauh lebih banyak daripada
\qty{0.40}{\percent} yang ada: \omterm{def:b2:magma-emplacement:magma}{magma} \omterm{def:b2:magma-emplacement:bodies}{sill} tidak melakukan eksolusi air.

\textbf{16.} CO$_2$ jauh kurang larut daripada air: gas itu menjadi jenuh dan membentuk gelembung
jauh di awal kenaikan, sehingga gas pertama sebuah basal sebagian besar CO$_2$.

\textbf{17.} Per kilogram, $n = 4.0/18.015 = \qty{0.222}{mol}$, $V = 0.222\times8.314\times1423/
\num{e5} = \qty{0.0263}{m^3}$; lelehan menempati $0.996/2710 = \qty{3.68e-4}{m^3}$;
fraksinya $0.0263/(0.0263 + 0.000368) = 0.986$, yaitu \qty{98.6}{\percent}.

\textbf{18.} $(1150 + 100)/2 = \qty{625}{\celsius}$.

\textbf{19.} Setelah 10 tahun, $2\sqrt{\kappa t} = 2\sqrt{\num{0.7e-6}\times\num{3.156e8}} =
\qty{29.7}{m}$. Titik itu berada pada $x = \qty{60}{m}$ dari tengah: $T = 100 +
525\,[\operatorname{erf}(-10/29.7) + \operatorname{erf}(110/29.7)] = 100 + 525\,(-0.365 +
1.000) = \qty{433}{\celsius}$.

\textbf{20.} Di tengah, $T = 100 + 1050\operatorname{erf}(50/2\sqrt{\kappa t})$: setelah
10 tahun $1050\operatorname{erf}(1.68) = 1050\times0.983$, $T = \qty{1132}{\celsius}$; setelah
50 tahun, $2\sqrt{\kappa t} = \qty{66.5}{m}$, $1050\operatorname{erf}(0.752) = 1050\times0.713$,
$T = \qty{848}{\celsius}$.

\textbf{21.} Pada $x = \qty{130}{m}$: \qty{294}{\celsius} setelah 300 tahun,
\qty{295}{\celsius} setelah 363 tahun (maksimum), \qty{295}{\celsius} setelah 420
tahun. Tanpa kalor laten, batuan \qty{80}{m} di atas \omterm{def:b2:magma-emplacement:bodies}{sill} tidak pernah benar-benar
mencapai \qty{300}{\celsius}.

\textbf{22.} $\num{9e8}\times1.772/(\num{2.5e6}\times1050) = 0.608$.

\textbf{23.} $\lambda = 0.64$: $e^{-0.4096}/(0.64\times1.6346) = 0.6639/1.0461 = 0.635$;
$0.65$: $0.6554/1.0673 = 0.614$; $0.66$: $0.6469/1.0886 = 0.594$. Nilai $0.608$
terletak di antara $0.65$ dan $0.66$, dekat dengan $0.65$: $\lambda = 0.65$.

\textbf{24.} $100 + 1050/(1 + 0.642) = 100 + 640 = \qty{740}{\celsius}$. Kalor
laten yang dilepaskan di front pembekuan menjaga sisi \omterm{def:b2:magma-emplacement:magma}{magma} tetap pada
\qty{1150}{\celsius} lebih lama, sehingga kontak menerima lebih banyak panas.

\textbf{25.} $(50)^2/\num{0.7e-6} = \qty{3.57e9}{s}$, yaitu 113 tahun.

\textbf{26.} $t = (100)^2/(16\times0.65^2\times\num{0.7e-6}) = \num{e4}/\num{4.73e-6} =
\qty{2.11e9}{s}$, sehingga $\num{2.11e9}/\num{3.156e7}$: \omterm{def:b2:magma-emplacement:bodies}{sill} itu padat setelah
\textbf{\qty[text-series-to-math=true]{67}{} tahun}. Waktu ini lebih pendek daripada
$(b/2)^2/\kappa$ dengan faktor $1/(4\lambda^2) = 0.59$: skala waktu itu hanya menetapkan
orde besarannya, dan kedua dinding \omterm{def:b2:magma-emplacement:bodies}{sill} menarik panas ke batuan induk yang dingin
sekaligus, cukup cepat untuk mengeluarkan kalor laten juga.
\end{solution}
